\section{Istanze, Varianti e Confronto Benchmark}

\subsection{Caratterizzazione delle Istanze di Test}
Per sottoporre a una valutazione empirica rigorosa gli algoritmi risolutivi e le relative varianti, è stato predisposto un insieme di istanze di prova strutturato attorno a dati realistici provenienti dal mondo del baseball professionistico (\textit{MLB}). La scelta di attingere a dati reali risponde all'obiettivo di testare il comportamento delle procedure non solo su distribuzioni sintetiche o puramente teoriche, ma su valori contrattuali e indicatori di rendimento che riflettono fedelmente le dinamiche decisionali del problema originario.

\subsubsection{Composizione del Dataset e Ruoli Difensivi}
Nel rispetto della formulazione matematica precedentemente definita, ciascuna istanza considera la copertura dei canonici $N = 9$ ruoli difensivi del baseball:
\begin{enumerate}
	\item \textit{Starting Pitcher} (lanciatore partente);
	\item \textit{Relief Pitcher} (lanciatore di rilievo);
	\item \textit{Catcher} (ricevitore);
	\item \textit{First Base} (prima base);
	\item \textit{Second Base} (seconda base);
	\item \textit{Third Base} (terza base);
	\item \textit{Shortstop} (interbase);
	\item \textit{Outfielder} (esterno);
	\item \textit{Designated Hitter} (battitore designato).
\end{enumerate}
Per ciascuno dei 9 ruoli è disponibile un contingente omogeneo di $P = 30$ candidati svincolati, per un totale complessivo di $N \times P = 270$ atleti presenti nel bacino di mercato.

Ciascun atleta $k \in \{1, \dots, P\}$ all'interno del ruolo $i \in \{1, \dots, N\}$ è formalmente caratterizzato da una tripla di attributi:
\begin{equation}
	p_{i, k} = \big( \mathit{name}_{i, k}, \; \mathit{cost}_{i, k}, \; v_{i, k} \big)
\end{equation}
dove $\mathit{cost}_{i, k} \in \mathbb{N}^+$ rappresenta il compenso contrattuale lordo (con valori compresi tra poche centinaia di migliaia di dollari fino a decine di milioni di dollari), mentre $v_{i, k} = \text{WAR} \ge 0$ rappresenta il contributo marginale del giocatore alle vittorie stagionali attese della squadra rispetto al livello di rimpiazzo.

\subsubsection{Generazione e Parametrizzazione degli Scenari}
Per analizzare sistematicamente la crescita del tempo di esecuzione e del consumo di memoria primaria, sono stati generati 10 scenari incrementali di test, denominati da \texttt{scenario\_11\_N9\_B3500\_P30} a \texttt{scenario\_20\_N9\_B8000\_P30}.

Nella progettazione della campagna sperimentale si è scelto di mantenere fisso il numero di atleti disponibili (confermati a $N = 9$ e $P = 30$ per tutte le istanze) e di far variare esclusivamente il budget nominale $X$. Il budget parte da un valore base di $\$350\,000\,000$ nello scenario 11 e cresce con incremento lineare e uniforme di $\$50\,000\,000$ per ciascuno scenario successivo, attestandosi infine a $\$800\,000\,000$ nello scenario 20.

Di conseguenza, il budget discreto associato varia progressivamente nell'intervallo:
\begin{equation}
	B \in \{3500, 4000, 4500, 5000, 5500, 6000, 6500, 7000, 7500, 8000\}
\end{equation}
La complessità di un'istanza nell'approccio dinamico è determinata dal prodotto delle tre dimensioni cardinali del problema, formalizzata dalla metrica:
\begin{equation}
	\text{instance\_size} = N \cdot B \cdot P
\end{equation}
Tale grandezza rappresenta l'esatto numero di transizioni valutate dal ciclo più interno della formulazione bottom-up classica. All'interno della suite di test, la dimensione dell'istanza cresce linearmente con pendenza costante da $945\,000$ fino a oltre $2\,160\,000$ operazioni elementari.

In Tabella~\ref{tab:scenarios_overview} vengono dettagliati i parametri operativi di ciascuno dei 10 scenari generati.

\begin{table}[!htbp]
	\centering
	\small
	\caption{Caratteristiche e Parametri delle Istanze di Test Generate}
	\label{tab:scenarios_overview}
	\begin{tabular}{lcccccc}
		\toprule
		\textbf{Identificativo Scenario} & $N$ & $P$ & \textbf{Budget Nominale ($X$)} & $B$ & \textbf{Dimensione ($N \cdot B \cdot P$)} \\
		\midrule
		\texttt{scenario\_11\_N9\_B3500\_P30} & 9 & 30 & $\$350\,000\,000$ & 3500 & $945\,000$ \\
		\texttt{scenario\_12\_N9\_B4000\_P30} & 9 & 30 & $\$400\,000\,000$ & 4000 & $1\,080\,000$ \\
		\texttt{scenario\_13\_N9\_B4500\_P30} & 9 & 30 & $\$450\,000\,000$ & 4500 & $1\,215\,000$ \\
		\texttt{scenario\_14\_N9\_B5000\_P30} & 9 & 30 & $\$500\,000\,000$ & 5000 & $1\,350\,000$ \\
		\texttt{scenario\_15\_N9\_B5500\_P30} & 9 & 30 & $\$550\,000\,000$ & 5500 & $1\,485\,000$ \\
		\texttt{scenario\_16\_N9\_B6000\_P30} & 9 & 30 & $\$600\,000\,000$ & 6000 & $1\,620\,000$ \\
		\texttt{scenario\_17\_N9\_B6500\_P30} & 9 & 30 & $\$650\,000\,000$ & 6500 & $1\,755\,000$ \\
		\texttt{scenario\_18\_N9\_B7000\_P30} & 9 & 30 & $\$700\,000\,000$ & 7000 & $1\,890\,000$ \\
		\texttt{scenario\_19\_N9\_B7500\_P30} & 9 & 30 & $\$750\,000\,000$ & 7500 & $2\,025\,000$ \\
		\texttt{scenario\_20\_N9\_B8000\_P30} & 9 & 30 & $\$800\,000\,000$ & 8000 & $2\,160\,000$ \\
		\bottomrule
	\end{tabular}
\end{table}

\subsubsection{Formato di Serializzazione e Struttura dei Dati}
Per garantire modularità, manutenibilità e riproducibilità degli esperimenti, ciascuna istanza è archiviata all'interno della cartella \texttt{scenarios/} come file autonomo in formato JSON. La separazione fisica tra la definizione dei dati di input e il codice degli algoritmi consente di modificare, estendere o aggiungere scenari senza dover alterare i moduli software.

La struttura di ciascun file JSON è articolata nei seguenti campi:
\begin{itemize}
	\item \texttt{scenario\_name}: stringa univoca identificativa dello scenario;
	\item \texttt{budget}: valore numerico intero del budget finanziario nominale $X$;
	\item \texttt{cost\_unit}: fattore di scala per la discretizzazione intera ($100\,000$);
	\item \texttt{roles}: vettore di $N$ oggetti, ciascuno rappresentante un ruolo difensivo e contraddistinto dal nome del ruolo (\texttt{role\_name}) e dalla lista dei $P$ atleti candidati (\texttt{candidates}). Ciascun elemento di tale lista riporta il nome del giocatore (\texttt{name}), il costo contrattuale lordo (\texttt{cost}) e il punteggio (\texttt{war}).
\end{itemize}

All'avvio dell'applicazione, il modulo \texttt{loader.py} esegue il parsing del documento JSON convertendolo negli oggetti immutabili \texttt{Scenario}, \texttt{Role} e \texttt{Player} definiti in \texttt{models.py}, validando preventivamente l'integrità dei parametri prima dell'invocazione degli algoritmi risolutivi.

\subsection{Varianti Algoritmiche Implementate}
Con l'obiettivo di ottimizzare le prestazioni computazionali ed estendere le dimensioni delle istanze elaborabili senza incorrere in saturazione di memoria, sono state sviluppate e sottoposte a sperimentazione diverse varianti dell'algoritmo di calcolo del valore ottimo, confrontate poi sistematicamente con la formulazione canonica di riferimento.

\subsubsection{Variante Bottom-Up con Sola Matrice \texorpdfstring{$m$}{m} (Senza Tabella \texorpdfstring{$s$}{s})}
L'algoritmo canonico (Algoritmo~\ref{alg:compute_optimal_war}) compie l'allocazione congiunta di due tabelle bidimensionali: la matrice dei valori $m$ e la matrice delle decisioni $s$. Come prima variante mirata al contenimento dell'occupazione spaziale, la procedura viene privata dell'allocazione e dell'aggiornamento della matrice ausiliaria $s$, calcolando e restituendo esclusivamente la matrice dei valori $m$.

Il codice Python implementato è il seguente:

\begin{lstlisting}
def compute_optimal_war_without_s(N, X, players, cost_unit=100000):
    B = floor(X / cost_unit)
    m = [[0.0 for _ in range(B + 1)] for _ in range(N + 1)]

    for i in range(1, N + 1):
        m[i][:] = m[i - 1][:]

        for p in players[i - 1]:
            c = p.cost // cost_unit
            war_p = p.war

            for b in range(c, B + 1):
                war = m[i - 1][b - c] + war_p
                if war > m[i][b]:
                    m[i][b] = war

    optimal_war = m[N][B]
    return optimal_war, m
\end{lstlisting}

Sotto il profilo asintotico, la complessità temporale rimane invariata a $\Theta(N \cdot B \cdot P)$, ma l'occupazione di memoria viene ridotta da $2(N + 1)(B + 1)$ a $(N + 1)(B + 1)$ celle, conseguendo un risparmio secco del $50\%$ dello spazio tabellare.

\subsubsection{Variante Salvaspazio a Vettore Singolo 1D}
Nell'algoritmo classico, per calcolare il valore ottimo si usa una tabella con $N+1$ righe e $B+1$ colonne, occupando uno spazio in memoria pari a $\Theta(N \cdot B)$. Tuttavia, quando si calcolano i valori per il ruolo $i$, servono solo ed esclusivamente i valori del ruolo precedente $i-1$. Non serve quindi conservare tutta la storia delle righe precedenti: basta usare un unico vettore di dimensione $B+1$.

Il punto chiave di questa variante è la direzione del ciclo sul budget, che deve andare obbligatoriamente \textbf{all'indietro, da $B$ verso $0$}:
\begin{itemize}
	\item Se scorressimo il budget in avanti (da $0$ a $B$), quando aggiorniamo una cella con budget $b$ usando il valore alla posizione $b - c$, rischieremmo di leggere un valore già modificato per lo stesso ruolo $i$. Questo permetterebbe all'algoritmo di scegliere più giocatori per lo stesso ruolo (o lo stesso giocatore più volte), violando la regola che impone al massimo un acquisto per ruolo.
	\item Scorrendo invece il budget all'indietro (da $B$ a $0$), quando andiamo a leggere la cella $b - c$, quella posizione non è ancora stata toccata per il ruolo $i$ e contiene ancora il valore del ruolo precedente $i-1$. In questo modo siamo sicuri che ogni giocatore venga considerato al massimo una volta per il suo ruolo.
\end{itemize}

Il codice Python implementato è il seguente:

\begin{lstlisting}
def compute_optimal_war_with_one_vector(N, X, players, cost_unit=100000):
    B = floor(X / cost_unit)
    solution = [0.0 for _ in range(B + 1)]

    for i in range(1, N + 1):
        for b in range(B, -1, -1):
            best_war = solution[b]

            for p in players[i - 1]:
                c = p.cost // cost_unit
                war_p = p.war

                if b >= c:
                    war = solution[b - c] + war_p
                    if war > best_war:
                        best_war = war

            solution[b] = best_war

    optimal_war = solution[B]
    return optimal_war
\end{lstlisting}

Il tempo di calcolo rimane invariato ($\Theta(N \cdot B \cdot P)$), perché il numero di controlli da fare è lo stesso. Il vantaggio sta nella memoria, che scende da $\Theta(N \cdot B)$ a \textbf{$\Theta(B)$}, diventando del tutto indipendente dal numero di ruoli $N$.

L'unico compromesso di questo approccio è che, sovrascrivendo i valori riga dopo riga nello stesso vettore, alla fine dell'esecuzione si perde la traccia delle scelte fatte e non è possibile ricostruire quali specifici giocatori compongono la rosa ottima.

\subsubsection{Variante Ricorsiva Top-Down con Memoizzazione}
L'approccio \textit{top-down} implementa la relazione di ricorrenza di Bellman mediante chiamate ricorsive a partire dallo stato obiettivo $(N, B)$ verso i casi base, memorizzando i risultati parziali in una matrice ausiliaria \texttt{memo} di dimensione $(N + 1) \times (B + 1)$, pre-inizializzata con valori sentinella negativi per distinguere gli stati non ancora calcolati.

Il codice Python implementato è il seguente:

\begin{lstlisting}
def compute_optimal_war_with_recursion(N, X, players, cost_unit=100000):
    B = floor(X / cost_unit)
    memo = [[float(-b - 1) for b in range(B + 1)] for _ in range(N + 1)]

    def dp(i, b):
        if i == 0:
            return 0.0

        if memo[i][b] >= 0:
            return memo[i][b]

        best_war = dp(i - 1, b)

        for p in players[i - 1]:
            c = p.cost // cost_unit
            if c <= b:
                candidate_war = dp(i - 1, b - c) + p.war
                if candidate_war > best_war:
                    best_war = candidate_war

        memo[i][b] = best_war
        return best_war

    optimal_war = dp(N, B)
    return optimal_war
\end{lstlisting}

Sotto il profilo teorico, il caso peggiore richiede la valutazione dell'intera tabella, mantenendo una complessità temporale $O(N \cdot B \cdot P)$ e spaziale $\Theta(N \cdot B)$ (a cui si sommano gli $O(N)$ record di attivazione generati sulla pila di chiamate ricorsive). Tuttavia, in scenari pratici in cui i costi dei giocatori presentano combinazioni discrete e intervallate, molti stati $(i, b)$ intermedi non risultano mai raggiungibili partendo da $(N, B)$. La variante top-down calcola esclusivamente gli stati strettamente necessari alla soluzione, consentendo in contesti sparsi un marcato abbattimento dei tempi effettivi di calcolo.

\subsubsection{Ottimizzazione mediante Filtro di Pareto (Potatura dei Candidati Dominati)}
Un'ulteriore ottimizzazione consiste nella riduzione a priori del bacino dei candidati disponibili per ciascun ruolo tramite l'estrazione della \textit{frontiera di Pareto}.
Nel problema in esame, per ciascun ruolo si seleziona al più un singolo atleta. Considerati due candidati $p_1, p_2$ appartenenti al medesimo ruolo $i$, diciamo che $p_1$ \textbf{domina} $p_2$ se:
\begin{equation}
	c_{i, p_1} \le c_{i, p_2} \quad \wedge \quad v_{i, p_1} \ge v_{i, p_2}
\end{equation}
con almeno una disuguaglianza stretta. Un candidato dominato richiede una spesa maggiore (o uguale) per offrire un rendimento inferiore (o uguale); esso non potrà pertanto mai far parte di alcuna configurazione ottima e può essere scartato senza alterare l'ottimalità della soluzione finale.

L'algoritmo di pre-filtraggio opera in due fasi:
\begin{enumerate}
	\item Per ogni ruolo, elimina gli atleti con $v \le 0$ o con costo eccedente il budget totale $B$, e ordina i restanti candidati per costo crescente (e WAR decrescente in caso di pareggio) in tempo $O(P \log P)$.
	\item Esegue una scansione lineare in tempo $O(P)$ mantenendo il massimo WAR osservato fino a quel punto, conservando solo i candidati che presentano un WAR strettamente superiore al massimo corrente.
\end{enumerate}

\begin{lstlisting}
def filter_players(players, cost_unit=100000, total_budget=None):
    valid_players = []
    for players_list in players:
        role_valid = []
        for p in players_list:
            c = p.cost // cost_unit
            if p.war > 0.0 and (total_budget is None or c <= total_budget):
                role_valid.append(p)
        role_valid.sort(key=lambda p: (p.cost, -p.war))
        valid_players.append(role_valid)

    filtered_players = []
    for role_candidates in valid_players:
        role_valid = []
        max_war = 0.0
        for p in role_candidates:
            if p.war > max_war:
                role_valid.append(p)
                max_war = p.war
        filtered_players.append(role_valid)
    return filtered_players
\end{lstlisting}

Il costo computazionale preliminare è $O(N \cdot P \log P)$, trascurabile rispetto all'esecuzione della programmazione dinamica. L'effetto è quello di ridurre il numero effettivo di candidati da $P$ a un sottoinsieme $P' \ll P$, velocizzando proporzionalmente il ciclo interno di tutte le varianti bottom-up che non necessitano della corrispondenza biunivoca degli indici nella matrice $s$.
